\section{Implementation Readiness and Systems Alignment}
\label{sec:systems}

\begin{table}[t]
\centering
\small
\caption{Runtimes across modalities comparing DcCA with state-of-the-art baselines. N/A indicates unsupported stages. All times are in milliseconds (ms): preprocessing and indexing are measured over the dataset, while search (Step~2.a) and set selection (Step~2.b) are reported per query.}
\resizebox{\linewidth}{!}{
\begin{tabular}{l r l l l r r r r}
\specialrule{0.12em}{0pt}{0pt}
Method & Dim & Feature & Engine & Retrieval & Prep. & Index & Search & Set \\
\specialrule{0.12em}{0pt}{0pt}

\multicolumn{9}{c}{\textbf{Tabular} ($N$=10000)} \\
\midrule
Baseline & 20 & Raw & DuckDB & Linear Scan & 0.00 & 23.62 & 2.66 & N/A \\
KD-Tree  & 20 & Raw & SciPy & KD-Tree Search & 0.00 & 3.57 & 0.98 & N/A \\
DcCA     & 64 & Table2Vec & FAISS & k-NN Search & 14.03 & 0.87 & 3.91 & 3.03 \\

\specialrule{0.12em}{0pt}{0pt}
\multicolumn{9}{c}{\textbf{Temporal} ($N$=500)} \\
\midrule
Baseline & 50 & Raw & tslearn & DTW Scan & 0.00 & 0.00 & 2377.50 & N/A \\
SAX      & 10 & Symbolic & iSAX & Hamming Search & 0.93 & 0.00 & 0.49 & N/A \\
DcCA     & 128 & Time2Vec & FAISS & k-NN Search & 2.05 & 0.06 & 1.28 & 2.51 \\

\specialrule{0.12em}{0pt}{0pt}
\multicolumn{9}{c}{\textbf{Text} ($N$=500)} \\
\midrule
Baseline & 512 & Keyword & DuckDB-FTS & Keyword Match & 0.00 & 2.75 & 1.12 & N/A \\
BM25     & 512 & BM25 & Whoosh & BM25 Ranking & 0.00 & 29.25 & 14.67 & N/A \\
DcCA     & 384 & SBERT & FAISS & k-NN Search & 253.01 & 0.09 & 1.57 & 2.82 \\

\specialrule{0.12em}{0pt}{0pt}
\multicolumn{9}{c}{\textbf{Image} ($N$=100)} \\
\midrule
Baseline & 150528 & Raw & DuckDB-Blob & Linear Scan & 0.00 & 0.00 & 103.75 & N/A \\
dHash    & 64 & dHash & FastBit & Hamming Search & 0.09 & 0.00 & 0.09 & N/A \\
DcCA     & 512 & ResNet18 & FAISS & k-NN Search & 1921.81 & 0.74 & 3.69 & 3.20 \\

\specialrule{0.12em}{0pt}{0pt}
\multicolumn{9}{c}{\textbf{Graph} ($N$=2000)} \\
\midrule
Baseline & N/A & Adjacency & Neo4j & BFS Search & 0.00 & 6.30 & 0.17 & N/A \\
DcCA     & 128 & GCN & Neo4j-VEC & k-NN Search & 3.73 & 0.14 & 1.52 & 2.60 \\

\specialrule{0.12em}{0pt}{0pt}
\multicolumn{9}{c}{\textbf{MultiModal: Tabular + Text + Image} ($N$=100)} \\
\midrule
DcCA     & 512 & CLIP & FAISS & k-NN Search & 2265.25 & 0.14 & 2.32 & 2.82 \\

\specialrule{0.12em}{0pt}{0pt}
\end{tabular}
}
\label{tab:multimodal_comparison}
\end{table}


Recent progress in modern data systems indicates that counterfactual querying is becoming increasingly aligned with mainstream data infrastructure, rather than being confined to the application layer. Vector search capabilities are now embedded in platforms such as PostgreSQL (pgvector) \cite{pgvector}, Neo4j \cite{neo4jVector}, Snowflake \cite{snowflakeCortex}, and Oracle \cite{oracleHybridIndex,oracleSearch}, reflecting a broader shift toward integrating similarity search directly into database engines and query execution pipelines \cite{johnson2019billion}. These systems already offer essential primitives, including scalable retrieval, structured filtering, hybrid lexical--semantic matching \cite{yates2021pretrained,guu2020retrieval}, and reranking \cite{khattab2020colbert}.

Nevertheless, the presence of these primitives does not equate to native support for counterfactual querying. Current systems typically decompose retrieval, filtering, and reranking into independent pipeline stages \cite{croft2010search,miutra2018introduction}, a design widely adopted in both database and information retrieval settings. In contrast, counterfactual querying requires these components to be jointly optimized under explicit constraints. The notion of counterfactual validity introduces conditions that fundamentally reshape the search space, such as opposite-outcome predicates and feasibility constraints, thereby turning the problem into a constrained, objective-driven retrieval task.

This gap is evident in existing system behavior. For instance, pgvector performs filtering after approximate index traversal \cite{pgvector}, which can shrink the set of valid candidates and necessitate repeated scans. Similar issues arise in approximate nearest neighbor search, where post-filtering under structured constraints leads to reduced recall \cite{dong2011efficient,malkov2018efficient}. These observations point to a central challenge: counterfactual constraints are not merely post-processing steps but actively redefine the effective search space, demanding tighter coupling among indexing, filtering, and ranking mechanisms.

From a data management standpoint, the core limitation lies not in scalability, but in the lack of native abstractions for counterfactual querying. Existing systems do not treat opposite-outcome predicates, feasibility-aware constraints, or set-level objectives such as diversity, coverage, robustness, and fairness as first-class query constructs, despite increasing interest in richer query semantics and learned retrieval objectives \cite{pachera2025whatif}.

This perspective also differentiates our approach from prior database frameworks grounded in explicit causal modeling or intervention semantics \cite{pachera2025whatif}. Instead, we adopt a data-centric formulation, where counterfactuals are retrieved directly from observed data under structured constraints, aligning with trends in retrieval-augmented reasoning \cite{lewis2020retrieval} and example-based explanations \cite{schultheis2023overview}.

Taken together, these insights suggest that counterfactual querying should be promoted to a first-class abstraction in data management systems, enabling joint specification and optimization of retrieval, constraints, and objectives on top of existing system primitives.